% ============================================================================
%  Relatorio Tecnico - Trabalho 4 de Programacao Paralela (Fases Paralelas,
%  MPI puro x hibrido MPI/OpenMP)
%  Formato exigido (mesmo do T1-T3):
%    pag. 1  -> coluna DUPLA, margens 1cm, fonte 10, somente texto
%    pag. 2  -> coluna SIMPLES: tabelas + graficos (anexo A)
%    pag. 3+ -> coluna SIMPLES: codigo fonte
%  Compilar a partir de relatorio/ (usa ../resultados/ e ../ para os fontes):
%    tectonic relatorio.tex    ou    pdflatex relatorio.tex
%  \pend{...} marca em VERMELHO o que depende dos numeros do cluster.
%  Antes de entregar: grep pend relatorio.tex deve voltar so a definicao.
%  Os numeros saem de resultados/tabelas.txt (python3 analisar.py).
% ============================================================================
\documentclass[10pt,a4paper,twocolumn]{article}

\usepackage{iftex}
\ifPDFTeX
  \usepackage[utf8]{inputenc}
  \usepackage[T1]{fontenc}
\fi
\usepackage[brazil]{babel}
\usepackage[margin=1cm]{geometry}
\usepackage{graphicx}
\usepackage{booktabs}
\usepackage{listings}
\usepackage{xcolor}
\usepackage{amsmath}

\setlength{\parindent}{0pt}
\setlength{\parskip}{0pt}
\setlength{\columnsep}{6mm}
\pagestyle{empty}
\raggedbottom
\makeatletter
\renewcommand\section{\@startsection{section}{1}{0pt}%
  {0.6ex plus .2ex minus .2ex}{0.2ex}{\normalsize\bfseries}}
\makeatother

\newcommand{\pend}[1]{\textcolor{red}{\textbf{[#1]}}}

% Inclui o grafico/tabela se ja existir; senao deixa um aviso no lugar.
\newcommand{\grafico}[2][\linewidth]{\IfFileExists{../resultados/#2.pdf}%
  {\includegraphics[width=#1]{../resultados/#2.pdf}}%
  {\pend{falta \detokenize{#2}.pdf: rodar python3 analisar.py}}}
\newcommand{\tabela}[1]{\IfFileExists{../resultados/#1.tex}%
  {\input{../resultados/#1.tex}}{\pend{falta \detokenize{#1}.tex: rodar python3 analisar.py}}}

\lstdefinestyle{codigoC}{
  language=C, basicstyle=\ttfamily\scriptsize, numbers=left,
  numberstyle=\tiny\color{gray}, keywordstyle=\color{blue!70!black},
  commentstyle=\itshape\color{green!40!black}, stringstyle=\color{red!60!black},
  showstringspaces=false, breaklines=true, tabsize=4,
  columns=fullflexible, lineskip=-1pt,
  literate={á}{{\'a}}1 {ã}{{\~a}}1 {ç}{{\c{c}}}1 {ê}{{\^e}}1 {í}{{\'i}}1
           {ó}{{\'o}}1 {ú}{{\'u}}1 {é}{{\'e}}1 {õ}{{\~o}}1 {à}{{\`a}}1
           {×}{{$\times$}}1 {→}{{$\rightarrow$}}1 {−}{{$-$}}1 {°}{{$^\circ$}}1,
  morekeywords={MPI_Init_thread,MPI_Sendrecv,MPI_Allreduce,MPI_Reduce,
                MPI_Barrier,MPI_Abort,MPI_PROC_NULL,omp_get_max_threads}
}

\begin{document}

% ============================ CABECALHO =====================================
\twocolumn[
{\centering
  {\large\bfseries Difusão de Calor (Jacobi) no Modelo de Fases Paralelas:\\ MPI Puro × Híbrido MPI/OpenMP}\\[2pt]
  {\small Trabalho 4 --- Programação Paralela --- Grupo 20 (conta LAD \texttt{cp32020})}\\[1pt]
  {\small Gabriel Coelho \quad Gabriel Cherubini}\\[4pt]
}]

% ============================ TEXTO (pag. 1) =================================
\section{Problema e adequação ao modelo}
Uma placa de $N \times N$ pontos com a borda de cima a 100\,°C e as outras
a 0\,°C. O método de Jacobi repete, para todo ponto interno,
$u'_{i,j} = \tfrac14(u_{i-1,j} + u_{i+1,j} + u_{i,j-1} + u_{i,j+1})$ até a
maior variação $\delta = \max|u' - u|$ ficar abaixo de uma tolerância. O
problema encaixa nas três fases do modelo: cada processo atualiza a sua faixa
de linhas (\emph{processamento local}), todos decidem juntos se $\delta$ já
está abaixo da tolerância (\emph{verificação coletiva}) e, se não, trocam com
os vizinhos de cima e de baixo as linhas de borda de que a próxima iteração
precisa (\emph{troca para convergência}). O trabalho local é um laço de
pontos independentes, que as threads dividem sem sincronização.

\section{Implementação: diretivas MPI e OpenMP}
Um fonte (\texttt{jacobi.c}) gera as três versões: \texttt{gcc} dá o
sequencial; \texttt{mpicc -DPARALELO}, o MPI puro; e o mesmo com
\texttt{-fopenmp}, o híbrido (l.~= linha do \texttt{jacobi.c}, Anexo B). As $N$ linhas são divididas em faixas
contíguas de $\lfloor N/np \rfloor$ ou $\lceil N/np \rceil$ linhas, geradas no
próprio processo (l.~201--216), cada uma com uma linha fantasma (\emph{halo}) acima e outra
abaixo. Cada iteração (laço das l.~234--258): (1) \texttt{passo()} (l.~104--122)
atualiza a faixa e devolve o $\delta$ local; (2) \texttt{MPI\_Allreduce(MPI\_MAX)} dá a todos o
$\delta$ global (l.~245), e todos param na mesma iteração (l.~251). Ele faz o papel dos $np$
\texttt{MPI\_Bcast} do pseudocódigo em uma coletiva de $\log_2 np$ etapas;
(3) dois \texttt{MPI\_Sendrecv} trocam a primeira linha com o vizinho de cima
e a última com o de baixo (\texttt{MPI\_PROC\_NULL} nas pontas, l.~220--221), sem
impasse (l.~130--141).

\textbf{Versão híbrida.} O modelo não muda: muda quem faz a fase 1. Um
processo por nó (ou por socket) recebe uma faixa 16 (ou 8) vezes maior, e
\texttt{\#pragma omp parallel for schedule(static) reduction(max:dmax)} (l.~108)
divide as linhas entre as threads; a redução do OpenMP combina o $\delta$ das
threads antes da redução do MPI. Só a thread principal chama MPI, fora das
regiões paralelas (\texttt{MPI\_Init\_thread} com
\texttt{MPI\_THREAD\_FUNNELED}, l.~169). A inicialização (l.~90--98) usa o mesmo laço e o
mesmo escalonamento, para que cada página de memória seja tocada primeiro pela
thread que vai usá-la (\emph{first touch}) e fique no banco de memória do
socket dela.

\section{Carga e metodologia}
$N = 8192$ (67\,M pontos, 2 matrizes de 0,5\,GB) e 500 iterações; a
tolerância ($10^{-6}$) é verificada a cada iteração, mas não é atingida, de
modo que todas as versões fazem o mesmo trabalho. Cada linha de borda tem
64\,KiB. No híbrido com 4 processos, cada faixa tem 2048 linhas, ou 128 por
thread com 16 threads: carga de sobra para dividir entre as threads.
O sequencial leva 138,9\,s. Atlântica, 4 nós exclusivos (2$\times$ Xeon E5520
por nó: 2 sockets $\times$ 4 núcleos $\times$ 2 HT), gcc 9.4.0 \texttt{-O2},
Open MPI 4.1.1. MPI puro e híbrido são comparados com o mesmo total de
núcleos e as mesmas CPUs: com 8, um rank (ou thread) por socket; com 16, dois;
com 32, todos os núcleos; com 64, também o HT. Processos e threads ficam
fixos (\texttt{--map-by ppr}, \texttt{OMP\_PLACES}/\texttt{OMP\_PROC\_BIND};
mapa conferido com \texttt{-m}). Tempo do laço principal, mediana de 3
(variação entre repetições $\le 2{,}3$\,\%);
$S = T_{seq}/T$ e $E = S/(\text{processos} \times \text{threads})$. O
resultado (iterações, $\delta$ bit a bit e soma das temperaturas) é igual ao
sequencial em todas as execuções.

\section{Escalabilidade forte}
Com 4 núcleos (1 por nó) as duas versões dão $S = 3{,}83$ ($E = 0{,}96$);
com 8 (1 por socket), 7,16 no MPI puro e 7,02 no híbrido. Daí em diante as
curvas achatam (Tab.~1, Fig.~1): com 16, 32 e 64 núcleos o MPI puro chega a
8,00, 8,67 e 8,51, e o híbrido com 1 processo por nó a 7,70, 7,51 e 7,47; a
eficiência cai a 0,13. O HT não ajuda (64 contra 32), e o híbrido com 1
processo por socket fica igual ao de 1 por nó (7,57 e 7,50). O híbrido com 1
thread repete o MPI com 4 processos (36,32 contra 36,26\,s): a medição é justa.

\section{Análise dos resultados}
\textbf{O gargalo é a memória, não o processador.} Cada ponto faz 5
operações e move $\approx$24\,B (lê a linha nova, escreve o resultado e a
leitura que precede a escrita). Um núcleo sozinho já faz 241\,M pontos/s,
$\approx$5,8\,GB/s. Com 4 núcleos por socket (32 no total) o cálculo leva
14,9\,s contra 16,5\,s com 1 por socket (8 no total): 4$\times$ os núcleos,
10\,\% a menos de tempo. O barramento de cada socket satura com cerca de um
núcleo, e o speed-up acompanha o número de controladores de memória em uso
(4 sockets $\to$ 3,8; 8 $\to$ 7,2), não o de núcleos. O HT divide o mesmo
barramento e por isso não muda nada.

\textbf{Menos mensagens não bastam.} O híbrido troca 6 mensagens por iteração
contra 126 do MPI com 64 processos, mas o volume que cruza a rede é o mesmo
(3 fronteiras entre nós, $2 \times 64$\,KiB cada), e as trocas dentro do nó
são cópias em memória compartilhada, baratas. O que pesa é quem espera a
rede. No híbrido, os processos 1 e 2 têm vizinhos em outros nós dos dois
lados e fazem as duas trocas pela rede em série: a troca de bordas custa em
média $\approx$4\,ms por iteração por processo (2,0\,s em 500 iterações, o
mesmo em todas as configurações $4 \times t$). No MPI puro, cada processo de fronteira tem só um vizinho em
outro nó. Como o \texttt{Allreduce} sincroniza todos a cada iteração, essa
espera entra no caminho crítico: verificação + bordas somam 1,1\,s no MPI
com 32 e 2,7\,s no híbrido $4 \times 8$ (Fig.~2). O híbrido ainda calcula
$\approx$5\,\% mais devagar (15,8 contra 14,9\,s), provavelmente pela barreira
implícita no fim de cada laço OpenMP, que espera a thread mais lenta. O
desbalanceamento entre processos é pequeno (cálculo máximo até 4,6\,\% acima
da média).

\section{Uso de IA}
% CONFERIR com o grupo antes de entregar: descrever exatamente o que foi feito
% pela ferramenta e o que foi feito por cada integrante (e se os slides da
% apresentacao tambem foram feitos com a ferramenta).
Usamos o \textbf{Claude Code} (Anthropic, modelo Claude Opus 5.5).
\emph{Código}: escreveu o \texttt{jacobi.c} (as três versões), os testes de
corretude, os scripts de medição e análise e o mapeamento de processos e
threads, que conferiu num simulador de topologia do Open MPI.
\emph{Relatório}: redigiu o texto e o Anexo A a partir dos números medidos.
O grupo escolheu o problema, rodou os experimentos
no cluster e revisou os resultados e o texto.

\section{Conclusão}
Neste problema e nesta máquina, a banda de memória limita as duas versões:
o speed-up para perto de 8,7 com os 8 sockets, e nem o HT nem o OpenMP mudam
isso. O MPI puro foi um pouco melhor (8,67 contra 7,51 com 32 núcleos)
porque cada processo de fronteira espera uma troca pela rede, e não duas.
O híbrido deve compensar com muito mais nós (o \texttt{Allreduce} e o número
de mensagens crescem com $np$), com um cálculo mais pesado por ponto, ou
esperando menos a rede: com \texttt{MPI\_Isend}/\texttt{MPI\_Irecv}, as duas
trocas dos processos do meio podem correr juntas, e uma thread dedicada à
comunicação sobreporia as bordas ao cálculo do interior.

% ============================ ANEXO (pag. 2) ================================
\onecolumn
\section*{\large Anexo A --- Tabelas e gráficos}
{\small
\begin{center}
\textbf{Tabela 1.} Escalabilidade forte ($N = 8192$, 500 iterações): tempo,
speed-up $S$, eficiência $E$ e fração do tempo médio por fase
(cálculo / verificação / troca de bordas).\\[2pt]
\tabela{tab_escala}

\vspace{6pt}
\grafico[0.9\linewidth]{grafico_escala}\\[2pt]
\textbf{Figura 1.} Escalabilidade forte de MPI puro e híbrido; tracejado = ideal;
faixa bege = região com HT (64 = 4 nós $\times$ 16 hardware threads).

\vspace{6pt}
\grafico[0.9\linewidth]{grafico_fases}\\[2pt]
\textbf{Figura 2.} Fração do tempo médio por processo em cada fase
($p \times t$ = processos $\times$ threads; H-nó = 1 processo por nó, H-sk = 1 por socket).
\end{center}
}

% ============================ CODIGO (pag. 3+) ==============================
\newpage
\section*{\large Anexo B --- Código}
\lstinputlisting[style=codigoC]{../jacobi.c}
\lstinputlisting[style=codigoC,language=make]{../Makefile}
\lstinputlisting[style=codigoC,language=bash]{../teste.sh}
\lstinputlisting[style=codigoC,language=bash]{../medir.sh}
\lstinputlisting[style=codigoC,language=bash]{../job.sh}
\lstinputlisting[style=codigoC,language=Python]{../analisar.py}

\end{document}
